\documentclass[10pt]{article}

\newcommand{\ResumeItemSep}{0.06em}
\newcommand{\ResumeTopSep}{0.10em}
\newcommand{\ResumeArrayStretch}{1.05}
\newcommand{\ResumeLineSpread}{1.00}
\newcommand{\ResumeSectionBefore}{0.24em}
\newcommand{\ResumeSectionAfter}{0.06em}

\usepackage{hyperref}
\usepackage{xcolor}
\usepackage{calc}
\usepackage{graphicx}
\usepackage{tikz}
\usepackage{fontspec}
\usepackage{fontawesome5}
\usepackage{titlesec}
\usepackage{enumitem}
\usepackage{fancybox}
\usepackage{xeCJK}

\hypersetup{hidelinks}
\urlstyle{same}

%%%%%%%%%%%%%%%%%%%%
% 设置
%%%%%%%%%%%%%%%%%%%%

\providecommand{\ResumeItemSep}{0.12em}
\providecommand{\ResumeTopSep}{0.18em}
\providecommand{\ResumeArrayStretch}{1.12}
\providecommand{\ResumeLineSpread}{1.08}
\providecommand{\ResumeSectionBefore}{0.55em}
\providecommand{\ResumeSectionAfter}{0.28em}

\setlength{\parindent}{0pt}
\setlength{\parskip}{0pt}
\pagenumbering{gobble}
\setlist[itemize]{
    leftmargin=1.35em,
    itemsep=\ResumeItemSep,
    topsep=\ResumeTopSep,
    parsep=0pt,
    partopsep=0pt
}
\setlist[enumerate]{leftmargin=*}
\renewcommand{\arraystretch}{\ResumeArrayStretch}
\linespread{\ResumeLineSpread}

\titleformat{\section}
  {\Large\bfseries\raggedright}
  {}{0em}
  {}
  [{\color{secondarycolor}\titlerule}]
\titlespacing*{\section}{0pt}{\ResumeSectionBefore}{\ResumeSectionAfter}

\usepackage[
    a4paper,
    left=1.1cm,
    right=1.1cm,
    top=0.75cm,
    bottom=0.85cm,
    nohead
]{geometry}

% 字体设置
% 说明：Noto Serif SC 没有原生斜体字型。将 CJK 的 italic 字型映射到常规字形
%（与过去“找不到斜体、默认替换为常规字形”的实际渲染一致），可消除
% “Font shape `TU/NotoSerifSC(0)/m/it' undefined” 告警；如需真正斜体观感，
% 可在此行追加 ItalicFeatures={FakeSlant=0.18}。
\setCJKmainfont[
    Path=fonts/,
    Extension=.otf,
    BoldFont=*-Bold,
    ItalicFont={NotoSerifSC},
]{NotoSerifSC}

% 自定义颜色
\definecolor{primarycolor}{RGB}{200, 60, 60}
\definecolor{secondarycolor}{RGB}{200, 68, 28}

\newlength{\iconwidth}
\setlength{\iconwidth}{1.5em}

\newcommand{\sectionicon}[2]{\makebox[\iconwidth][c]{\color{primarycolor}{#1}}\quad #2}

\newcommand{\projectheading}[3]{%
    \noindent\makebox[\textwidth][s]{%
        \makebox[0pt][l]{{\large \textbf{#1}}}%
        \hfill
        \makebox[0pt][c]{{\normalsize \textbf{#2}}}%
        \hfill
        \makebox[0pt][r]{#3}%
    }%
}

% 通用个人信息：修改这里即可同步到所有岗位模板。
\newcommand{\ResumeName}{无绝}
\newcommand{\ResumeBirthDate}{2002.10}
\newcommand{\ResumeHometown}{湖北咸宁}
\newcommand{\ResumeSchool}{重庆邮电大学}
\newcommand{\ResumeEnglishLevel}{CET-6}
\newcommand{\ResumeEmail}{250xxxxxxx@qq.com}
\newcommand{\ResumePhone}{198xxxxxxxx}
\newcommand{\ResumeHomepageLabel}{github.com/Dithob}
\newcommand{\ResumeHomepageUrl}{https://github.com/Dithob}
\newcommand{\ResumePhoto}{images/xjpic.jpg}
\newcommand{\ResumeHeaderImage}{images/foot.png}

% 通用教育背景：修改这里即可同步到所有岗位模板。
\newcommand{\ResumeGraduateSchool}{重庆邮电大学}
\newcommand{\ResumeGraduateMajor}{计算机技术}
\newcommand{\ResumeGraduateDegree}{硕士}
\newcommand{\ResumeGraduateDate}{2024.09--2027.06}
\newcommand{\ResumeGraduateAwards}{特等学业奖学金 $\times$ 1、二等学业奖学金 $\times$ 1}

\newcommand{\ResumeUndergraduateSchool}{重庆邮电大学}
\newcommand{\ResumeUndergraduateMajor}{计算机科学与技术}
\newcommand{\ResumeUndergraduateDegree}{本科}
\newcommand{\ResumeUndergraduateDate}{2020.09--2024.07}
\newcommand{\ResumeUndergraduateAwards}{三等学业奖学金；GPA: 3.28/4.0（专业前 25\%）}

\newcommand{\ResumeEducationBlock}{%
    {\textbf{\ResumeGraduateSchool}}，\ResumeGraduateMajor，\textit{\ResumeGraduateDegree} \hfill \ResumeGraduateDate
    \begin{itemize}
        \item \textbf{荣誉奖项}：\ResumeGraduateAwards
    \end{itemize}

    \vspace{0.2em}
    {\textbf{\ResumeUndergraduateSchool}}，\ResumeUndergraduateMajor，\textit{\ResumeUndergraduateDegree} \hfill \ResumeUndergraduateDate
    \begin{itemize}
        \item \textbf{荣誉奖项}：\ResumeUndergraduateAwards
    \end{itemize}
}


%%%%%%%%%%%%%%%%%%%%
% 文章内容
%%%%%%%%%%%%%%%%%%%%

\newcommand{\ResumeTargetRole}{前端开发}
\providecommand{\ResumeTargetRole}{求职方向待定}

\newcommand{\ResumeContact}{%
    \footnotesize
    \textcolor{white}{%
        \faEnvelope \quad 邮箱：\href{mailto:\ResumeEmail}{\ResumeEmail}
        \hspace{3.2em}
        \faPhone \quad 电话：\ResumePhone
        \hspace{3.2em}
        \faGithub \quad 主页：\href{\ResumeHomepageUrl}{\ResumeHomepageLabel}
    }%
}

\newcommand{\ResumeContactBar}{%
    \begin{tikzpicture}[remember picture, overlay]
        \node[anchor=north, inner sep=0pt](headercontacts) at (current page.north){%
            \includegraphics[width=\paperwidth]{\ResumeHeaderImage}
        };
        \node[anchor=center] at(headercontacts.center){\ResumeContact};
    \end{tikzpicture}
    \vspace{-1.15em}
}

\newcommand{\ResumeProfileSection}{%
    \section[个人信息]{\sectionicon{\faAddressCard}{个人信息}}
    \begin{tabular*}{\textwidth}{@{\extracolsep{\fill}}lll}
        \textbf{姓\qquad 名}：\ResumeName & \textbf{出生年月}：\ResumeBirthDate & \textbf{籍\qquad 贯}：\ResumeHometown \\
        \textbf{学\qquad 校}：\ResumeSchool & \textbf{英语水平}：\ResumeEnglishLevel & \textbf{求职方向}：\ResumeTargetRole
    \end{tabular*}
}

\newcommand{\ResumeEducationSection}{%
    \section[教育背景]{\sectionicon{\faGraduationCap}{教育背景}}
    \ResumeEducationBlock
}

\newcommand{\ResumeProfileAndEducation}{%
    %%%%%%%%%%%%%%%%%%%%
    % 个人信息与教育背景
    %%%%%%%%%%%%%%%%%%%%

    \begin{minipage}[t]{0.79\textwidth}
        \ResumeProfileSection

        \ResumeEducationSection
    \end{minipage}
    \hfill
    \begin{minipage}[t]{0.14\textwidth}
        \vspace{-0.35em}
        \setlength{\fboxsep}{0pt}
        % 双线框 (\doublebox) 会给照片额外加约 4pt 边框宽；图片若占满 \linewidth 会触发
        % “Overfull \hbox (3.99988pt too wide)” 告警，故图片宽度预留 4pt，使含框总宽恰好等于栏宽。
        \doublebox{\includegraphics[width=\dimexpr\linewidth-4pt\relax]{\ResumePhoto}}
    \end{minipage}
}


\begin{document}

    % 顶部联系人栏
    \ResumeContactBar
    \ResumeProfileAndEducation

    % \vspace{0.45em}

    %%%%%%%%%%%%%%%%%%%%
    % 专业技能
    %%%%%%%%%%%%%%%%%%%%

    \section[专业技能]{\makebox[\iconwidth][c]{\color{primarycolor}{\faWrench}}\quad 专业技能}
    \begin{itemize}
        \item \textbf{前端基础与框架}：熟悉 HTML5、CSS3、JavaScript、TypeScript，掌握 Vue3、Pinia，具备组件化开发、状态管理、路由组织与复杂交互实现经验；具备 Next.js 数据看板开发实践。
        \item \textbf{组件与可视化}：熟悉 Ant Design Vue、Element 等组件库，能够封装表单、列表、弹窗、搜索筛选等通用业务组件；熟悉 ECharts，可实现基础数据看板与可视化展示。
        \item \textbf{工程实践与协作}：掌握 Git、Docker、Postman、JMeter、FoxAPI 等工具，熟悉前后端接口联调、请求/响应数据排查、环境部署与基础性能测试流程。了解 SpringBoot、MyBatis 等后端框架，掌握 MySQL、SQLite、Redis 的基本使用，能够参与接口设计、数据库建模和前后端数据流打通。
        \item \textbf{AI 辅助研发}：熟悉 Claude Code、Codex 等 AI 编程工具，具备 Skill 编排、MCP 服务开发与集成经验，能够结合研发代码上下文开展自动化用例生成与公共方法复用。具备 SSE 流式对话、模型工具调用展示、实时预览与可视化编辑开发经验，使用 \texttt{EventSource}、\texttt{iframe} 与 \texttt{postMessage} 实现相关交互。
        % \item \textbf{AI 辅助开发}：熟悉 GitHub Copilot、Claude Code、Codex、Qoder 等 AI 编程工具，能够结合 AI 对话、SSE 流式响应和实时预览能力开发 AI 应用型前端项目。
        
        
        
    \end{itemize}

    % \vspace{0.45em}

    %%%%%%%%%%%%%%%%%%%%
    % 实习经历
    %%%%%%%%%%%%%%%%%%%%

    \section[实习经历]{\makebox[\iconwidth][c]{\color{primarycolor}{\faBriefcase}}\quad 实习经历}

    \projectheading{腾讯云与智慧产业研发公司}{前端开发实习生}{2026.06--2026.09}
    \begin{itemize}
        \item \textbf{地图数据看板开发}：基于 Next.js 与腾讯地图 JSAPI GL 开发 LBS 营销数据看板模板，负责地图可视化页面与业务数据展示，支撑业务分析结果交付。
        \item \textbf{Web 地图前端组件开发与维护}：基于 TypeScript 与腾讯地图 JSAPI GL 开发和维护 Web 地图组件，封装地图初始化、图层渲染与交互事件等通用逻辑，支撑业务数据的地图可视化展示；维护 GL API 与数据可视化示例，完善组件调用、数据展示与交互逻辑。
        \item \textbf{手机地图H5 }：参与「你用我赔H5」活动页面与业务交互开发，对接后端接口实现活动状态展示与跨天状态更新，接入端外分享及埋点上报；完善移动端布局与多端兼容性适配，支撑 H5 活动上线。
    \end{itemize}

    % \vspace{0.25em}

    % \projectheading{重庆啄木鸟网络科技有限公司}{前端开发实习生}{2025.06--2025.09}
    % \begin{itemize}
    %     \item \textbf{核心职责}：参与家电售后场景 AI 应用研发，负责需求分析、接口设计与前后端联调，参与智能产品识别、小啄 GPT 智能回复、Agent 文档解析与联网搜索等模块。
    %     \item \textbf{智能家电识别与自动入库系统}：基于 Dify 搭建“识别--联网补全--字段校验--向量去重--入库”工作流，并通过接口封装为前端提供结构化产品信息与补全、去重结果数据。
    % \end{itemize}

    % \vspace{0.45em}

    %%%%%%%%%%%%%%%%%%%%
    % 项目经历
    %%%%%%%%%%%%%%%%%%%%

    \section[项目经历]{\makebox[\iconwidth][c]{\color{primarycolor}{\faProjectDiagram}}\quad 项目经历}

    \projectheading{地图移动端 GUI 自动化质检 Agent}{参与框架建设}{2026.06--2026.09}
    \begin{itemize}
        \item \textbf{Agent 编排与用例生成}：构建 LLM Agent 质量效能体系，按职责编排 13 个 Skill 与 3 个 MCP 服务，打通「需求 → 自动化用例生成 → 知识沉淀」闭环；联动 TAPD/Figma 自动生成可执行 case.py，单用例生成约 10 分钟（手写需 2 小时），覆盖 17 个业务域、600+ 生成代码单元。
        % \item \textbf{Agent 编排与工具集成}：参与自动化质检框架建设，编排 13 个 Skill 与 3 个 MCP 服务，联动 TAPD、Figma 生成可执行自动化用例，单用例代码生成耗时约 10 分钟。
        \item \textbf{MCP 工具与多模态感知}：自研设备控制与线上诊断 MCP，并接入地图 MCP 解析坐标构造测试数据；融合四层视觉感知并链式降级，解决自绘控件无 resource-id 的定位问题；scrcpy 长连接投屏将截图耗时从约 1s 降至 50ms。
        \item \textbf{知识工程与语义检索}：设计页面/组件/坑点/操作四类节点知识图谱，基于 23 个研发仓库 UI 索引做同源静态推理，让 Agent 生成前先读研发代码。 AST 扫描 1500+ 私有 helper，结合 jieba + BM25 实现公共方法语义检索与经验复用。
        % \item \textbf{代码知识检索与复用}：结合研发仓库 UI 索引与 AST 扫描构建代码知识索引，使用 BM25 实现公共方法检索与经验复用，为 Agent 用例生成提供研发代码上下文。
    \end{itemize}

    % \vspace{0.18em}

    % \projectheading{云端智能图像协同系统}{前端开发}{2025.06--2025.08}
    % \begin{itemize}
    %     \item \textbf{项目目标}：面向公共图库、个人私有图库与企业团队共享场景，开发支持图片检索、多维分析、管理员审核与 AI 实时协同编辑的图像管理平台。
    %     \item \textbf{权限与任务状态}：基于 Vue Router 为路由节点配置 \texttt{meta.access} 字段，并结合 \texttt{beforeEach} 全局守卫实现页面级访问控制；对接后端 AI 图像处理模型，针对“AI 扩图”等高耗时任务采用“任务派发--定时轮询”机制追踪异步执行状态。
    %     \item \textbf{协同编辑与可视化}：基于 WebSocket 事件驱动架构开发多人协同编辑能力，引入编辑锁与冲突检测机制，降低并发编辑下的状态覆写风险；使用 ECharts 封装分组条形图、词云图等数据组件，并借助 Vue3 响应式能力动态更新图表配置。
    % \end{itemize}
    % ※ 2026-09-14 用户裁定注释停用，勿自动恢复。

    \vspace{0.18em}

    \projectheading{智能应用生成平台}{前端开发}{2025.09--2025.12}
    \begin{itemize}
        \item \textbf{AI 应用交互开发}：面向非技术用户，开发多轮 AI 对话、响应式网站预览与可视化修改功能，支持通过自然语言生成网站；基于 Ant Design Vue 封装应用卡片、信息弹窗等业务组件。
        \item \textbf{流式响应与过程反馈}：基于原生 \texttt{EventSource} 对接后端 SSE 接口，实现 AI 响应内容的增量展示，并动态呈现模型工具调用轨迹，向用户反馈网站生成过程。
        \item \textbf{预览通信与可视化编辑}：通过 \texttt{postMessage} 实现主应用与 \texttt{iframe} 预览页面之间的消息通信，同步预览与编辑状态；在预览页面内接入 DOM 事件监听脚本，实现页面元素点选与可视化修改。
    \end{itemize}
    % \begin{itemize}
    %     \item \textbf{项目目标}：面向非技术用户的 AI 应用生成平台，支持通过多轮 AI 对话生成网站，并提供实时响应式预览与可视化修改能力。
    %     \item \textbf{流式交互与预览通信}：采用原生 \texttt{EventSource} 消费后端 SSE 流式接口，实现 AI 对话打字机输出效果，并动态展示模型工具调用轨迹；使用 \texttt{postMessage} 打通主应用与 \texttt{iframe} 预览站点之间的跨域通信，支撑实时预览与编辑状态同步。
    %     \item \textbf{可视化编辑与组件复用}：通过向子 \texttt{iframe} 动态注入 DOM 监听脚本，实现“点击即修改”的可视化编辑能力；基于 Ant Design Vue 封装应用卡片、信息弹窗等业务组件，通过 \texttt{props}/\texttt{emits} 进行父子组件通信，提高页面开发复用率。
    % \end{itemize}

    % \vspace{0.25em}
    % \projectheading{校园二手交易平台}{项目负责人}{2025.03--2025.06}
    % \begin{itemize}
    %     \item \textbf{项目目标}：面向校园二手交易场景，开发支持商品发布、商品浏览、关键词搜索与交易流程的信息平台，覆盖用户、商品与交易等核心实体管理。
    %     \item \textbf{页面与交互实现}：负责商品列表、商品详情、商品发布和搜索结果等核心页面开发，围绕列表渲染、条件筛选、表单提交和状态提示组织前端交互逻辑；基于响应式布局适配不同屏幕尺寸。
    %     \item \textbf{模板复用与数据流打通}：使用 Django 模板引擎抽取页面公共结构，实现导航、商品卡片、列表区域等模块复用；参与后端 API 设计与 SQLite 实体关系建模，打通商品数据展示、提交、查询和交易状态流转。
    % \end{itemize}
    % ※ JD 切换库存：投「通用前端 / 全栈（Django）」方向时，与「医药问数」项目互换启用。

    \vspace{0.18em}

    % \projectheading{基于 RAG 的医药数据分析问数 Agent}{项目主要成员}{2025.10--2026.02}
    % \begin{itemize}
    %     \item \textbf{项目目标}：面向药品采购、零售、库存等全链路数据及 800 余张业务表，建设 LLM 数据分析问答助手，将业务人员的自然语言问题转化为可执行 SQL 与可读分析结论。
    %     \item \textbf{链路设计}：负责 LLM 应用核心链路设计，打通「问题理解 → Schema 检索 → SQL 生成 → 结果解释 → 人工反馈」业务流程；聚焦结果解读与人工反馈环节的落地，降低业务人员直接理解复杂数据库结构的门槛。
    %     \item \textbf{结果可视化与 AI 报表}：负责问数结果的前端可视化呈现与 AI 报表输出，将查询结果转化为可视化图表与结构化报表；针对采购、零售、库存等业务场景设计数据看板与可视化大屏，支撑业务人员直观解读分析结论。
    %     \item \textbf{评估与自修正闭环}：引入 SQL 预执行与自纠错机制，捕获 MySQL Error 后回传模型重写；构建约 200 条 Golden Dataset，结合 RAGAS 指标、人工反馈 Bad Case 回流评估检索与 SQL 效果，执行成功率由约 70\% 提升至 90\%+，项目获重庆市人工智能大赛市级银奖。
    % \end{itemize}
    
    % \vspace{0.18em}
    
    \projectheading{医药智能问数与数据可视化平台}{项目主要成员}{2025.10--2026.02}
    \begin{itemize}
        \item \textbf{前端页面与问数交互}：基于 Vue3 + TypeScript + Ant Design Vue 开发自然语言问数、结果展示与人工反馈页面，组织问题输入、查询状态、可视化图表与 AI 分析结论，支持业务人员完成「提问 → 分析 → 反馈」交互流程。
    
        \item \textbf{ECharts 图表组件封装}：围绕采购、零售、库存场景，封装柱状图、折线图等可复用图表组件，通过 \texttt{props} 接收业务数据与展示配置，结合 \texttt{watch} 与 \texttt{setOption} 实现数据更新；处理图表容器尺寸变化与实例销毁，支持看板响应式展示。
    
        \item \textbf{动态数据适配与报表展示}：对接 800 余张业务表的 RAG 问数链路，封装查询结果的字段映射与数据转换逻辑，将维度、指标适配为图表配置和表格数据，统一呈现趋势图表与结构化报表；完善加载、空数据及查询异常状态展示。
    
        \item \textbf{AI 链路与质量评估}：参与 Schema 检索、SQL 生成与结果解释链路建设，引入 SQL 预执行和错误反馈重写机制；构建含约 200 条样本的评估集，结合页面人工反馈与 Bad Case 回流，推动 SQL 执行成功率由约 70\% 提升至 90\%+。
    \end{itemize}

    % \vspace{0.45em}

    %%%%%%%%%%%%%%%%%%%%
    % 奖项证书
    %%%%%%%%%%%%%%%%%%%%

    \section[奖项证书]{\makebox[\iconwidth][c]{\color{primarycolor}{\faTrophy}}\quad 奖项证书}
    \begin{itemize}
        % \item \textbf{证书}：CET-6
        \item \textbf{竞赛}：2025 “互联网+”大学生创新创业大赛银奖；2025 首届重庆市 AI 大模型创新应用大赛省级三等奖；2023 第十四届蓝桥杯 C/C++ 组省赛三等奖；2023 重庆市高校数据库应用竞赛校级二等奖
        \item \textbf{论文}：《LHAF: A Lightweight Hierarchical Adaptive Framework for LLM-based Multimodal Emotion Recognition in Conversations》（SCI I 区在投）
    \end{itemize}

\end{document}
